\subsection{模型四的建立与求解：固定轨迹任务优化}\label{sec:model-q4}

\subsubsection{模型四的建立}\label{sec:q4-establishment}

\noindent\textbf{（1）轨迹状态、目标距离与拍照方向角。}

问题四不改变问题三所得轨迹，只在固定轨迹上选择任务和时刻。设机器人状态为
\(\mathbf p(t)=(x(t),y(t))^{\mathsf T}\)、速度模 \(v(t)\)、加速度模
\(a(t)\)，目标 \(j\) 的坐标为 \(\mathbf g_j\)，则距离为

\begin{equation}
d_j(t)=\|\mathbf p(t)-\mathbf g_j\|_2.
\label{eq:q4-distance}
\end{equation}

射击任务在执行前连续 1.5 s 内须满足

\begin{equation}
5\le d_j(t)\le30,\qquad v(t)\le2,\qquad a(t)\le1.5;
\label{eq:q4-shooting-constraints}
\end{equation}

每个射击目标至多执行一次，单次命中概率为 0.85。拍照任务在执行前连续 0.5
s 内须满足

\begin{equation}
10\le d_j(t)\le40,\qquad v(t)\le1.5,\qquad a(t)\le1.5.
\label{eq:q4-photo-constraints}
\end{equation}

机器人相对拍照目标的方向角定义为

\begin{equation}
\theta_j(t)=\operatorname{atan2}(y(t)-y_j,x(t)-x_j)\pmod {360^\circ},
\label{eq:q4-photo-angle}
\end{equation}

同一目标两次拍照的圆周角距离不得小于
\(60^\circ\)。题面未明示单机器人能否并行准备，也未规定同一射击目标是否可重复计数。为使``完成目标''具有明确工程含义，本文作两条建模假设：机器人只有一套任务执行资源，任意两项已选任务的准备/执行闭区间不得相交，并额外留出
0.01 s
数值间隔；每个射击目标至多计一次完成。结果总数还受模板九行容量限制。这两条属于模型假设而非题面原文，后续约束与最优性结论均以此为边界。

\noindent\textbf{（2）连续可行窗口与候选任务。}

先在正式 10 Hz
状态上，对每个可能执行时刻向前检查完整准备窗的全部采样点。对距离下界、距离上界、速度和加速度分别归一化，候选
\(c\) 的最小裕度定义为

\begin{equation}
m_c=\min_{t\in[t_c-L_c,t_c]}
\left\{
\frac{d-d_{\min}}{d_{\max}-d_{\min}},
\frac{d_{\max}-d}{d_{\max}-d_{\min}},
\frac{v_{\max}-v}{v_{\max}},
\frac{a_{\max}-a}{a_{\max}}
\right\}.
\label{eq:q4-margin}
\end{equation}

在每个连续可行段内保留定时间隔点、局部高裕度点；拍照任务另按 30°
方向角箱保留高裕度点。随后用 PCHIP 连续状态在 0.01 s
网格重新检查完整准备区间，剔除任何负裕度候选。共得到 1255
个候选，其中射击 595 个、拍照 660 个。

\noindent\textbf{（3）资源互斥、射击唯一性与词典序MILP。}

令 \(x_c\in\{0,1\}\) 表示是否选择候选 \(c\)。对时间冲突候选对
\((c,d)\in\mathcal T\)、同一射击目标候选集
\(\mathcal S_j\)，以及同一拍照目标且角差不足 60° 的候选对
\((c,d)\in\mathcal A\)，约束为

\begin{equation}
x_c+x_d\le1\quad((c,d)\in\mathcal T\cup\mathcal A),
\qquad
\sum_{c\in\mathcal S_j}x_c\le1,
\qquad
\sum_cx_c\le9.
\label{eq:q4-logical-constraints}
\end{equation}

第一级目标为

\begin{equation}
\max\sum_cx_c.
\label{eq:q4-milp-stage1}
\end{equation}

固定最优任务数 \(N^*\) 后，引入最小裕度 \(z\)，以充分大的 \(M\) 施加

\begin{equation}
z\le m_c+M(1-x_c),\qquad \max z.
\label{eq:q4-milp-stage2}
\end{equation}

最后固定 \(N^*\) 和
\(z^*\)，最大化总裕度；对可同时选择的同目标拍照对增加 AND
辅助变量及按圆周角差归一化的多样性奖励，射击的 0.85
期望命中仅作为很小的同级择优项。三个阶段均由 HiGHS 返回最优状态且 MIP
gap 为
0。求解器的线性优化实现背景参见文献\cite{huangfu2018}，但本文的最优性结论只依据实际求解状态、零 gap 与独立复核。由于九行是不可突破的题面输出上界，找到九项可行任务也证明了任务数量已达全局上界；裕度和多样性的最优性针对经连续验证的压缩候选集。

\subsubsection{模型四的求解}\label{sec:q4-solution}

求解先识别连续可行窗口，再以0.01 s网格细化并压缩候选；随后依次最大化任务数量、最小安全裕度和总裕度/拍照多样性。三个阶段均固定前一阶段最优值，故次级目标不会牺牲任务数量。最终对候选窗口作0.001 s连续复核，并在时间保留两位小数后再次核验。

候选可行性分布见图\ref{fig:q4-candidates}。候选集中于轨迹经过目标邻域且速度、加速度同时满足完整准备窗的区段，这一分布来自物理约束，而不是优化器任意删选。

\begin{figure}[htbp]
  \centering
  \includegraphics[width=.90\textwidth]{q4/candidate_feasibility_map.pdf}
  \caption{模型四经连续准备窗复核后的1255个候选任务分布}
  \label{fig:q4-candidates}
\end{figure}

\subsubsection{算法实现}\label{sec:q4-algorithm}

\begin{pseudocodeblock}{算法4\quad 固定轨迹上的任务调度}
{\heiti 输入：}问题三10 Hz状态、模型内位置标准差、目标坐标及任务约束。\par
{\heiti 输出：}九项任务的目标编号、任务类型、准备时刻与执行时刻。\par
\begin{enumerate}[label={\heiti 步骤~\arabic*:},leftmargin=5.5em,itemsep=1pt,topsep=3pt]
  \item 插值轨迹状态，逐目标检查射击或拍照的完整准备区间，识别连续可行窗口；
  \item 在每个窗口内保留定时间隔点、局部高裕度点及拍照方向角代表点，并在0.01 s网格上复核；
  \item 建立时间冲突、角度冲突、射击唯一性和九行容量约束，依次求解任务数、最小裕度和综合质量三级MILP；
  \item 若任一级未返回最优状态或MIP gap不为0，则终止并报告失败；否则进入连续复核；
  \item 在0.001 s网格上检查所选任务，时刻保留两位小数后再次检查全部准备窗与冲突约束；
  \item 将通过复核的九行结果写入工作簿副本并返回日程，原始模板保持不变。
\end{enumerate}
\end{pseudocodeblock}

\subsubsection{结果分析}\label{sec:q4-results}

正式方案见表\ref{tab:q4-schedule}。输出时刻按模板保留两位小数，并在舍入后再次进行 0.01
s 全准备窗复核。

\begin{table}[htbp]
\centering
\caption{问题四正式九项任务日程}
\label{tab:q4-schedule}
\begin{tabular}{rccrr}
\toprule
序号 & 目标编号 & 任务 & 开始准备时刻/s & 任务执行时刻/s \\
\midrule
1 & S05 & 射击 & 450.31 & 451.81 \\
2 & P04 & 拍照 & 454.61 & 455.11 \\
3 & P09 & 拍照 & 455.21 & 455.71 \\
4 & P12 & 拍照 & 456.01 & 456.51 \\
5 & S04 & 射击 & 464.21 & 465.71 \\
6 & P16 & 拍照 & 466.61 & 467.11 \\
7 & P03 & 拍照 & 493.21 & 493.71 \\
8 & S11 & 射击 & 509.31 & 510.81 \\
9 & S10 & 射击 & 525.51 & 527.01 \\
\bottomrule
\end{tabular}
\end{table}

方案含 4 次射击、5 次拍照，射击期望命中数为

\begin{equation}
E(N_{\mathrm{hit}})=4\times0.85=3.40.
\label{eq:q4-expected-hits}
\end{equation}

五次拍照分属五个不同目标，因此同目标角差约束自动满足。正式方案最小归一化约束裕度为
0.4241；最早可行贪心基线虽同样完成九项，但最小裕度仅 0.1010，MILP
方案约为其 4.20
倍。全部任务在准备区间内的距离、速度和加速度均已通过连续复核。

轨迹、目标位置和所选任务见图\ref{fig:q4-layout}，日程时间线见图\ref{fig:q4-timeline}。任务集中在450.31--527.01 s是固定轨迹与目标几何共同作用的结果；所有准备/执行闭区间互不重叠，最紧相邻任务仍保留0.10 s间隔。

\subsubsection{模型检验与稳健性}\label{sec:q4-validation}

设置三档压力情景，每档使用固定随机种子生成 500 次扰动：中等情景采用 0.5
倍问题三模型位置标准差和 5\% 速度/加速度相对扰动，标准情景采用 1.0
倍位置标准差和 5\% 状态扰动，严重情景采用 1.5 倍位置标准差和 10\%
状态扰动。九项任务逐项和整套方案在三档情景中的可行率均为
100\%。该结果只说明方案在所定义扰动集下具有较大裕度，不被解释为真实任务成功概率；射击命中仍单独按题设
0.85 计算。

最终将九行结果写入 \texttt{result.xlsx} 副本的 B2:E10。A 列序号、H:L
红色说明和范例、其余单元格值及样式均经语义快照核对保持不变；原始模板未被覆盖。
约束裕度见图\ref{fig:q4-margins}。
全部1255个候选的裕度经验分布见图\ref{fig:q4-margin-cdf}。

\begin{figure}[htbp]
  \centering
  \includegraphics[width=.88\textwidth]{q4/trajectory_targets_schedule.pdf}
  \caption{模型四固定轨迹、目标点与九项所选任务}
  \label{fig:q4-layout}
\end{figure}
\begin{figure}[htbp]
  \centering
  \includegraphics[width=.88\textwidth]{q4/optimized_schedule_timeline.pdf}
  \caption{模型四九项任务的准备与执行时间线}
  \label{fig:q4-timeline}
\end{figure}
\begin{figure}[htbp]
  \centering
  \includegraphics[width=.90\textwidth]{q4/constraint_margins.pdf}
  \caption{模型四九项任务的归一化完整准备窗约束裕度}
  \label{fig:q4-margins}
\end{figure}
\begin{figure}[htbp]
  \centering
  \includegraphics[width=.78\textwidth]{paper/candidate_margin_cdf.pdf}
  \caption{模型四射击与拍照候选裕度的经验累积分布}
  \label{fig:q4-margin-cdf}
\end{figure}
